\chapter{82}



Unten empfing Louis ein dominant zuckriger Schokoladeduft. Der
erneute Szenenwechsel war kaum zu schaffen. Er hing dem Gespräch mit
Valentina nach, als wären sie noch immer verbunden. Er stellte sich vor, wie Bastiano an der Tür stand und ihr Gespräch
mithörte. Und je länger Louis nachdachte, desto weniger verstand
er, was Bastiano eigentlich wollte. Was er mit Giorgias Sturz überhaupt
zu tun hatte. Warum er Valentina bedrängte. Und dann fielen ihm
Valentinas Worte wieder ein. Sie hatte von Spuren, von Fingerabdrücken
gesprochen. Davon, dass Bastiano verhört worden war und sein Vater ihn
rausgeholt hatte. Woraus denn eigentlich?
Manuela und Ava standen mit dem Rücken zu ihm am grossen Tisch. Joh sass
ihnen gegenüber. Beide Hände hatte er über die Augen gelegt. Diesmal
schaffte es Louis nur schwer, seine Rolle zu spielen. 
Ava drehte sich zu ihm um.

«Augen zu, Ciro.»

Er liess sich von ihr am Arm wie eine Puppe zu einem Stuhl neben Joh
ziehen.

«Also,» begann Manuela, «jeder kann zwei Mal raten.»

Joh begann.

«Eismischung mit Schokosauce?»

Avas Nein kam siegessicher.

«Schokocrème?»

«Nein.»

Das Spiel erinnerte Louis an Kindergeburtstage. Doch selbst wenn es ihn albern anmutete und er den Eindruck hatte, das leise Surren verbunden mit dem unverkennbaren Schokoladenduft lasse keine Zweifel offen, sah er keinen andern Weg, als mitzutun.

«Schokobananen?»

«Nein.»

«Ihr kommt eh nicht drauf. Augen auf.»

Zu viert starrten sie auf die Tischmitte. Umrahmt von mehreren kleinen
Schälchen mit aufgespiessten Beeren stand ein kleines, silberfarbenes
Edelstahlkonstrukt. Langsam floss Schokolade über die Tropfteller nach
unten in eine Auffangschale.
Joh schnippte mit den Fingern.

«Ein Schokoladenbrunnen, ich glaub’s nicht.»

«Genau,» Manuela stemmte die Hände in die Hüften und sah aus wie eine
Matrone, «unsere Überraschung des Tages. Und was haben die Herren zu
bieten?»

Joh rollte die Augen.
Ava legte den Arm um Manuela.

«Ach, komm, Manuela, legen wir los,» und zu Joh und Louis gerichtet,
«überlegt euch was, bis wir fertig sind.»

Sie tunkten die Früchte in die Schokoladensauce, drehten sie darin herum
und steckten sie in den Mund. Für einen Moment war es ruhig. Fast
andächtig still.
Joh machte immer wieder ein überraschtes Gesicht.

«Sehr gut.»

Dann erzählte er, wie er den Zaun mit Louis erstellt hatte. Es klang
wie ein Rapport.
Schliesslich reckte sich Ava vor und fixierte Louis.

«Ich wüsste was, Ciro. Ich würde gerne einen Song von dir hören.»

Manuela stimmte ihr vergnügt zu.

«Ja, ja, das wäre toll.»

Auch das noch. Selbst wenn Louis solche Abende
kannte. Gesellige Runden endeten oft mit seinem Gitarrenspiel. Für
gewöhnlich liebte er diese Rolle. Jetzt aber war er wirklich zu aufgekratzt und
ahnungslos, was die drei hören wollten. Sie blickten ihn erwartungsvoll
an. Louis’ Gehirn ratterte. Wie würde er den Abend am leichtesten
überstehen? Eigentlich war es klar. Nichts viel ihm leichter, als in
seiner Musik unterzutauchen.

«Gut, aber nur unter der Bedingung, dass ihr Vorschläge macht. Wenn
nichts davon in meinem Repertoire ist, lassen wir es, okay?»

Seine Idee kam gut an.
Joh zeigte auf Ava. Mit einer Handbewegung lud er sie ein, die Wunschrunde
zu starten.

«Ladys first.»

«Hm, schwierig.»

Sie sah überrumpelt aus und senkte den Blick. Louis hatte den Eindruck,
sie ringe mit sich. Sie schaute zu Louis hoch. Nach drei
fehlgeschlagenen Musikwünschen hob sie den Zeigefinger.

«Aber Jovanotti kennst du doch bestimmt, nicht?»

«Ja, den schon.»

«Eines seiner Canzoni war unser Lieblingslied, also Daniels und meins. Seit Daniels Tod hab ich es mir nicht mehr angehört.»

Manuela berührte Avas Hand und wandte sich an Louis.

«Weisst du, Ciro, ihr Freund ist bei einem Motorradunfall verstorben.»

«Oh, das tut mir sehr leid.»

Louis erinnerte sich. Joh hatte ihn unter vier Augen informiert. Er
war erleichtert, dass Ava die Tragödie nun selbst angesprochen hatte.

«Was würdest du denn gerne hören?»

«Also, \emph{«Chiaro di Luna»} war so etwas wie unser Paar-Song,» sie
zögerte, «der Text passt, passte perfekt zu uns. Vor Daniels Tod waren wir in Italien an
einem Jovanotti-Konzert, es war unglaublich.»

Avas Blick verriet ein tief empfundenes Erlebnis. Ihre Augen schimmerten. Manuela streichelte ihr übers Haar.

\qrblock{Jovanotti \emph{«Chiaro di luna»}}{media/QR-Codes/043_Chiaro_di_Luna_Lorenzo_Jovannotti_Spotify.png}

Louis schluckte.

«Ja, ich kenne den Song, musikalisch und inhaltlich gross. Nur, den habe
ich nie gespielt und darum auch den Text nicht intus,» er runzelte
die Stirn, «aber ich spiele zum Beispiel \emph{«À Te»} von ihm. Kennst
du es?»

«Ja, das kenne ich, das ist auch sehr schön.»

Manuela und Joh waren zu stillen Beobachtern geworden. Sie wirkten, als überliessen
sie Ava und Louis das Feld. Zwischendurch glaubte Louis, auf Johs Lippen
ein verhaltenes Lächeln zu entdecken. Was wohl in ihm vorging?

«Gut, den Song spiel’ ich gerne für euch. \emph{«À te»} ist sehr speziell. Es wird gemunkelt, 
dass Jovanotti das Stück als Liebeserklärung an die Musik schrieb, nicht etwa an eine geliebte Frau,
wie man denken möchte; das hätten die Leute gern. Ich fühl mich mit
Jovanotti verbunden. Mich gäb’ s ohne Musik jedenfalls nicht.»

Ava sank zusammen.

«Aha.»

Sie suchte Louis’ Blick und sah enttäuscht aus. 

«Wie auch immer, wunderschön ist der Song ohnehin.»

Als Louis mit seiner Gitarre zurückkam, sassen die drei in den
tiefen Kissen des Sofas. Er setzte sich vor ihnen auf einen Stuhl.
Gleichzeitig fühlte er Bastiano neben sich. Als hätte der sich an ihm
festgebissen wie ein Tier. Und wenn «À te» auch gerade
weniger als nichts mit seiner Stimmung zu tun hatte, war er froh, sich für eine Songlänge in die Musik absetzen zu können. Er führte sein verletztes Bein in eine schmerzfreie Position. Dann
stimmte er die Gitarre. 

«Jovanotti begleitet den Song am Piano. Das mache ich für gewöhnlich
auch, aber jetzt hört ihr ihn mal mit Gitarrenbegleitung, und das
Orchester dazu stellt ihr euch einfach vor, hm?»

Manuela und Ava nickten.

«Toll, und das nächste Mal mit Piano,» Joh strahlte, «Manuela, es wird
endgültig Zeit, dass wir das Klavier runterholen.»

Manuelas Blick verriet Ungeduld.

«Ja, ja, gute Idee, aber jetzt freue ich mich, wenn es los geht.»

Sie stützte ihren Kopf auf den Händen ab. Und alle drei sahen bereit
aus.

\qrblock{Jovanotti \emph{«À Te»}}{media/QR-Codes/044_A_te_Lorenzo_Jovannotti_Spotify.png}

Avas Gesicht hatte von Takt zu Takt heller geleuchtet. Ihre Augen
ähnelten schmalen Blinkern. Es war nicht auszumachen, ob die Darbietung oder Louis selbst sie verzaubert, oder ob der Song sie in Zeiten mit Daniel zurückversetzt hatte.
Manuela winkelte ihre Beine an.

«Deine Stimme ist, wie ich schon sagte, Ciro, aussergewöhnlich - »

Joh fiel ihr ins Wort.

«Ja, sehe ich auch so,» er begann zu lachen, «wie sagt man so schön, das
Gesamtpaket stimmt, smarter Typ, tolle Stimme, virtuoser Musiker.»

Er schälte sich aus den Kissen.

«Wer mag ein Bier?»

Manuela hob die Hand. Louis nickte abwesend. Smarter Typ? Erschreckend verkehrt, wie Joh ihn wahrnahm. 
Ava sah aus, als käme sie von irgendwoher zurück.

«Danke, Ciro, der Song ist wunderschön, und dein dunkles Timbre in der
Stimme, wow, macht das Lied ganz besonders.»

Louis legte die Gitarre beiseite. Ava schlug die Augen nieder und schien wieder wegzudriften.
Selbst wenn seine Darbietung gelungen war, fühlte sich Louis fehl am Platz. Ciro verkam unaufhaltsam zu einer lächerlichen Figur. 
Johs Herumwerkeln in der Küche war ein Segen. Die Geräusche holten Louis’ 
Aufmerksamkeit für eine Weile aus der Enge, bevor er unweigerlich wieder bei Giorgia landete. 

Als Joh die Biere verteilte, drehten sich ihre Gespräche bald nur mehr um die Schafe. Johs spassig
dargebotene \emph{Flexa}-Geschichten schliesslich liessen den Abend
ausgelassen enden. Louis kostete es viel, Teil der Runde zu sein. 
